\documentclass[11pt]{amsart}
\usepackage{amsmath,amssymb,amsthm}
\usepackage[margin=1.1in]{geometry}
\usepackage{array,longtable,booktabs}
\usepackage{hyperref}
\emergencystretch=2em

\newtheorem{theorem}{Theorem}
\newtheorem{lemma}[theorem]{Lemma}
\theoremstyle{remark}
\newtheorem{remark}[theorem]{Remark}

\newcommand{\Inf}{\mathrm{Inf}}
\newcommand{\E}{\mathbb{E}}
\newcommand{\MAJ}{\mathrm{MAJ}}
\newcommand{\ii}{\mathrm{i}}

\title{A forced coordinate with nearly maximal erasure entropy}
\author{Yitzchak Shmalo}
\address{Einstein Institute of Mathematics, The Hebrew University of Jerusalem, Jerusalem, Israel}
\date{}

\begin{document}

\begin{abstract}
The erasure entropy $E_x(f)$ of a Boolean function at a coordinate $x$ is the conditional entropy of the event $x\in U$
given $U\setminus\{x\}$, where $U$ is drawn from the spectral distribution, and it never exceeds the binary entropy of
the influence $\Inf_x(f)$. We study an explicit monotone function on $2k+1$ variables in which one coordinate $x$ is
the only variable lying both in a shortest minterm and in a shortest maxterm. It is obtained from a majority function
with one hidden coordinate by replacing $(k-1)/2$ of the majority's $k$ inputs with the OR of two fresh variables and
the other $(k+1)/2$ with the AND of two. We prove that $E_x/\Inf_x=(1+o(1))\log_2(1/\Inf_x)=(1+o(1))(k-1)\log_2\frac43$, so the
ratio at the forced coordinate grows linearly in $k$ and is asymptotic to its largest possible value. The proof writes
the relevant biased Fourier coefficients of the majority as integrals over the unit circle, on which the integrand has
an explicit modulus with a single pair of peaks, and uses that the phase advances by $\arccos\frac13$, an irrational
multiple of $\pi$, from one class of coefficients to the next.
\end{abstract}

\subjclass[2020]{06E30, 42C10, 94A17}
\keywords{Fourier entropy-influence conjecture, Boolean functions, Fourier entropy, total influence}

\maketitle

\section{Introduction}\label{sec:intro}

For $f\colon\{0,1\}^n\to\{-1,1\}$ write $f=\sum_S\hat f(S)\chi_S$ and let $U$ be a random subset with
$\Pr[U=S]=\hat f(S)^2$. The Fourier entropy-influence conjecture of Friedgut and Kalai~\cite{FK96} asks for an absolute
$C$ with $H(U)\le C\,\E|U|$. It is known for symmetric functions and read-once decision trees~\cite{OWZ} and is
preserved under suitable compositions~\cite{OT13}. Weaker forms and bounds with extra factors are
in~\cite{CKLS,WWW,KKLMS,Han}. Li and Han~\cite{LH25} make the constant in front of $\E|U|$ in Han's bound
sharp. The conjecture was proved for every linear threshold function in~\cite{ShmLTF}. The largest known lower bound
on the best constant is $6.5218$~\cite{Num26}. Earlier bounds are due to O'Donnell and Tan~\cite{OT13}, Hod~\cite{Hod}
and Mosaic Intelligence~\cite{MI26a,MI26b}. The erasure entropy of $f$ at a coordinate $x$ is
\[
E_x(f)=H\bigl(\mathbf 1[x\in U]\bigm|U\setminus\{x\}\bigr),
\]
the conditional entropy of one coordinate given the others, in the sense of Verd\'u and Weissman~\cite{VW08}, and
with $\Inf_x(f)=\Pr[x\in U]$ the conjecture holds if and only if
\[
\sup_f\ \min_{x:\,\Inf_x(f)>0}\frac{E_x(f)}{\Inf_x(f)}
\]
is finite~\cite{Shm1}, which builds on the reduction of Gonz\'alez, MacManus and Pereyra~\cite{GMP} to a
statement about one coordinate; coordinates at which the ratio is bounded for some classes of functions are studied
in~\cite{Shm3}. Since conditioning does not increase entropy,
\begin{equation}\label{eq:upper}
\frac{E_x(f)}{\Inf_x(f)}\le\frac{h_2(\Inf_x(f))}{\Inf_x(f)}\le\log_2\frac{e}{\Inf_x(f)}
\end{equation}
for every $f$ and $x$ with $\Inf_x(f)>0$.

Which coordinate to take is the difficulty. In~\cite{Shm8} several natural selection rules are shown to fail on monotone
functions, among them the rule that takes, among the variables lying both in a shortest minterm and in a shortest
maxterm, one of smallest ratio. For that rule the counterexample is a family $F_{k,t}$ in which a coordinate $x$ is the
only variable of both kinds, and the ratio at $x$ is unbounded when $k$ and $t$ grow together. Here we treat the smallest case $t=1$, the simplest member of the family, and
find the exact rate of growth; the extension to every fixed $t$ and to every rational bias is in~\cite{Shm10}.

Fix an odd $k=2m+1$ and let $y^0\in\{0,1\}^k$ be the point with $y^0_j=1$ for $j\le m$ and $y^0_j=0$ for $j>m$,
and $y^1=y^0+e_{m+1}$. Let $M$ be the majority function $\MAJ_k$ with the point $y^1$ removed from its true set, and
\[
\tilde f_k(y,x)=M(y)\vee\bigl(x\wedge[y\ge y^0]\bigr).
\]
Let $F_k$ be the function of $2k+1$ variables obtained from $\tilde f_k$ by replacing $y_j$ with the OR of two fresh
variables for $j\le m$ and with the AND of two fresh variables for $j>m$. It is the function $F_{k,1}$
of~\cite[Section~5]{Shm8}; it is monotone, and $x$ is the only variable that lies both in a shortest minterm and in a
shortest maxterm of $F_k$~\cite[Section~5]{Shm8}.

\begin{theorem}\label{thm:main}
$\Inf_x(F_k)=(3/4)^{k-1}$ and, as $k\to\infty$ through odd values,
\[
\frac{E_x(F_k)}{\Inf_x(F_k)}=(1+o(1))\log_2\frac1{\Inf_x(F_k)}=(1+o(1))\,(k-1)\log_2\frac43 .
\]
\end{theorem}

The upper half of the theorem is~\eqref{eq:upper}. The exact values computed in Section~\ref{sec:num} approach the
bound slowly: the quotient of the ratio by $(k-1)\log_2\frac43$ is $0.80$ at $k=21$, $0.88$ at $k=81$ and $0.92$ at
$k=141$, increasing along the table, though not monotone in $k$ (its smallest value, $0.796$, is at $k=19$).

\section{The ratio as a double sum}\label{sec:sum}

All logarithms are to base $2$ unless written $\ln$. For $a,b\ge0$ let $\Phi(a,b)=a\log\frac{a+b}{a}+b\log\frac{a+b}{b}$,
with $\Phi(a,0)=\Phi(0,b)=0$, and $\Psi(r)=\Phi(r^2,1)$ for $r\ge0$. Then $\Psi\ge0$, and $\Psi(r)\ge\log r^2$ for
$r\ge1$, since both terms of $\Phi$ are nonnegative and $(r^2+1)/1\ge r^2$.

Put $P=\frac34$ and $q=\frac14$. The OR of two uniform bits is true with probability $P$ and the AND with probability
$q$. By the proof of~\cite[Proposition~12]{Shm8}, $E_x(F_k)$ and $\Inf_x(F_k)$ are the erasure entropy and the
influence of $x$ for the law $c(S)^2$, where $c(S)$ are the coefficients of $\tilde f_k$, in $\pm1$ values, in the
orthonormal basis $\prod_{j\in S}\phi_j(y_j)$ (times $2x-1$ when $x\in S$) of the product measure on $(y,x)$ with
$\Pr[y_j=1]=\mu_j$ and $x$ uniform; here $\mu_j=P$ for $j\le m$, $\mu_j=q$ for $j>m$, and
$\phi_j(y)=(y-\mu_j)/\sigma$ with $\sigma=\sqrt{Pq}=\sqrt3/4$ (the $\mu$-biased Fourier basis,
see~\cite[Chapter~8]{OD14}; the conjecture for biased product measures is formulated in~\cite{KMS12}).

With $+1$ for true, $\tilde f_k=\bar A+(2x-1)D$, where
\[
D=\mathbf 1[y=y^0]+\mathbf 1[y=y^1],\qquad \bar A=2N-1+\mathbf 1[y=y^0]-\mathbf 1[y=y^1],
\]
and $N=\mathbf 1[|y|\ge m+1]$. For $S\subseteq[k]$ write $u=|S\cap\{1,\dots,m\}|$ and $v=|S\cap\{m+2,\dots,k\}|$, and
let $\hat g(S)$ denote the coefficient of $g$ at $\prod_{j\in S}\phi_j$. Put
\[
X(z)=q+Pz,\qquad Y(z)=P+qz,\qquad H_{uv}(z)=(z-1)^{u+v}X(z)^{m-u}Y(z)^{m+1-v}.
\]

\begin{lemma}\label{lem:coef}
If $m+1\in S$ then $\hat D(S)=0$. If $m+1\notin S$ then
\[
\hat D(S)=(-1)^v\sigma^{u+v}P^{2m-u-v},\qquad
\hat{\bar A}(S)=2\sigma^{u+v}t_{uv}-\mathbf 1[S=\emptyset]+\tfrac12\hat D(S),
\]
where
\[
t_{uv}=\sum_{c\ge m+1}[z^c]H_{uv}(z).
\]
\end{lemma}

\begin{proof}
For a point $w$, $\widehat{\mathbf 1[y=w]}(S)=\Pr[y=w]\prod_{j\in S}\phi_j(w_j)$. For $j\le m$, $y^0_j=1$ and
$\phi_j(1)=q/\sigma=\sqrt{q/P}$; for $j\ge m+2$, $y^0_j=0$ and $\phi_j(0)=-q/\sigma=-\sqrt{q/P}$. The points $y^0$
and $y^1$ differ only at $m+1$, where $\mu_{m+1}=q$, so $\Pr[y=y^0]=P^{2m+1}$ and $\Pr[y=y^1]=P^{2m}q$. If $m+1\in S$
the two terms of $\hat D(S)$ carry the factors $P^{2m+1}\phi_{m+1}(0)=-P^{2m}\sqrt{Pq}$ and
$P^{2m}q\,\phi_{m+1}(1)=P^{2m}\sqrt{Pq}$, which cancel. Otherwise both points give the same product
$(-1)^v(q/P)^{(u+v)/2}$, so $\hat D(S)=P^{2m}(-1)^v(q/P)^{(u+v)/2}=(-1)^v\sigma^{u+v}P^{2m-u-v}$ and
$\widehat{\mathbf 1[y=y^0]}(S)-\widehat{\mathbf 1[y=y^1]}(S)=(P-q)\hat D(S)=\frac12\hat D(S)$. Finally
$\E[z^{y_j}\phi_j(y_j)]=\mu_jz(1-\mu_j)/\sigma-(1-\mu_j)\mu_j/\sigma=\sigma(z-1)$ and $\E[z^{y_j}]=1-\mu_j+\mu_jz$, so
by independence $\E[z^{|y|}\prod_{j\in S}\phi_j(y_j)]=\sigma^{u+v}H_{uv}(z)$ when $m+1\notin S$, and $\hat N(S)$ is the
sum of its coefficients at $z^c$, $c\ge m+1$.
\end{proof}

Let $\beta_u=\binom mu q^uP^{m-u}$ be the binomial law with parameters $m$ and $q$, and for $0\le u,v\le m$ put
\[
r_{uv}=\frac{\bigl|2t_{uv}-\mathbf 1[u=v=0]+\frac12(-1)^vP^{2m-u-v}\bigr|}{P^{2m-u-v}}.
\]

\begin{lemma}\label{lem:sum}
$\Inf_x(F_k)=P^{2m}=(3/4)^{k-1}$ and
\[
\frac{E_x(F_k)}{\Inf_x(F_k)}=\sum_{u,v=0}^m\beta_u\beta_v\,\Psi(r_{uv}).
\]
\end{lemma}

\begin{proof}
By the reduction above, $E_x=\sum_S\Phi(\hat{\bar A}(S)^2,\hat D(S)^2)$ and $\Inf_x=\sum_S\hat D(S)^2$, the sums over
$S\subseteq[k]$; the terms with $m+1\in S$ vanish. There are $\binom mu\binom mv$ sets with given $u,v$ and $m+1\notin
S$, and $\sum_{u,v}\binom mu\binom mv(Pq)^{u+v}P^{2(2m-u-v)}=(Pq+P^2)^{2m}=P^{2m}$. By homogeneity
$\Phi(\hat{\bar A}^2,\hat D^2)=\hat D^2\,\Psi(|\hat{\bar A}|/|\hat D|)$, the quotient $|\hat{\bar A}(S)|/|\hat D(S)|$ is
$r_{uv}$ by Lemma~\ref{lem:coef}, and $\binom mu\binom mv(Pq)^{u+v}P^{2(2m-u-v)}/P^{2m}=\beta_u\beta_v$.
\end{proof}

So $E_x/\Inf_x$ is the mean of $\Psi(r_{uv})$ when $u$ and $v$ are independent binomial variables with parameters $m$ and
$\frac14$. Most of that law sits near $u=v=m/4$, and the theorem says that there $\log r_{uv}^2$ is close to
$2m\log\frac43$ for almost all $(u,v)$. The difficulty is that $t_{uv}$ is a sum of terms of both signs; we show that it
oscillates, in $u$, like the cosine of an irrational rotation times an exponentially large envelope.

\section{The tail on the unit circle}\label{sec:circle}

\begin{lemma}\label{lem:integral}
If $u+v\ge1$ then $Q_{uv}=H_{uv}/(z-1)$ is a polynomial and
\[
t_{uv}=[z^m]\,Q_{uv}(z)=\frac1{2\pi}\int_{-\pi}^{\pi}Q_{uv}(e^{\ii\theta})\,e^{-\ii m\theta}\,d\theta .
\]
\end{lemma}

\begin{proof}
Write $Q_{uv}=\sum_cg_cz^c$. Then $[z^c]H_{uv}=g_{c-1}-g_c$, and the sum over $c\ge m+1$ telescopes to $g_m$ because
$Q_{uv}$ has degree one less than $H_{uv}$. The integral is Cauchy's formula on $|z|=1$; coefficient extraction of this kind, followed by a Laplace or
saddle-point estimate, is standard~\cite[Chapter~VIII]{FS09}.
\end{proof}

On the circle the three factors of $Q_{uv}$ have explicit moduli. For $z=e^{\ii\theta}$ and $w=\cos\theta$,
\[
|z-1|^2=2-2w,\qquad |X(z)|^2=|Y(z)|^2=q^2+P^2+2Pqw=\frac{5+3w}8 .
\]
Put $n_1=u+v-1$ and $n_2=2m+1-u-v$, so that $n_1+n_2=2m$. Then $|Q_{uv}(e^{\ii\theta})|=e^{\rho(\cos\theta)}$ with
\[
\rho(w)=\frac{n_1}2\ln(2-2w)+\frac{n_2}2\ln\frac{5+3w}8,\qquad
\rho''(w)=-\frac{n_1}{2(1-w)^2}-\frac{9n_2}{2(5+3w)^2}\le-\frac{9n_2}{128}.
\]
So $\rho$ is strictly concave on $[-1,1)$; its maximum is at $w_0=(3n_2-5n_1)/(6m)$ when this lies in $(-1,1)$, and a
short computation gives $2-2w_0=8n_1/(3m)$, $(5+3w_0)/8=n_2/(2m)$, hence
\begin{equation}\label{eq:peak}
\rho(w_0)=\frac{n_1}2\ln\frac{8n_1}{3m}+\frac{n_2}2\ln\frac{n_2}{2m}.
\end{equation}
At $u=v=m/4$ one has $n_1\approx m/2$, $w_0\approx\frac13$ and $e^{\rho(w_0)}\approx(\sqrt3/2)^m$.

For $\lambda\ge1$ let $W_\lambda$ be the set of $(u,v)$ with $|u-m/4|\le\lambda\sqrt m$ and $|v-m/4|\le\lambda\sqrt m$.
In what follows $c,C$ denote positive constants that may depend on $\lambda$ but not on $m$ or on $(u,v)\in W_\lambda$,
and $m$ is large enough for each statement. On $W_\lambda$, $n_1=m/2+O(\sqrt m)$ and $w_0=\frac13+O(m^{-1/2})$; put
$\theta_0=\arccos w_0$, $\theta^*=\arccos\frac13$ and $z_0=e^{\ii\theta_0}$.

\begin{lemma}\label{lem:exponent}
On $W_\lambda$, $\rho(w_0)+(n_2-1)\ln\frac43\ge(m-1)\ln\frac43$.
\end{lemma}

\begin{proof}
Regard the left side, by~\eqref{eq:peak}, as a function $G(n_1)$ of a real variable $n_1\in(0,2m)$, with $n_2=2m-n_1$.
Then $G'(n_1)=\frac12\ln\frac{16n_1}{3n_2}-\ln\frac43$, which vanishes at $n_1=m/2$, and
$G''(n_1)=\frac1{2n_1}+\frac1{2n_2}>0$. So $G$ is convex with its minimum at $m/2$, and
$G(m/2)=\frac m4\ln\frac43+\frac{3m}4\ln\frac34+(\frac{3m}2-1)\ln\frac43=(m-1)\ln\frac43$.
\end{proof}

Fix $(u,v)\in W_\lambda$ and let $\Lambda(\theta)=\ln Q_{uv}(e^{\ii\theta})-\ii m\theta$, for complex $\theta$ with
$|\theta-\theta^*|\le\delta$, where $\delta>0$ is so small that $e^{\ii\theta}$ stays away from $1$, $-\frac13$ and $-3$,
the zeros of $z-1$, $X$ and $Y$; the logarithm is the branch continuous in this disc. Since $\Lambda$ is a sum of at most
$2m+1$ logarithms of functions analytic and nonvanishing there, plus a linear term, $|\Lambda''|$ and $|\Lambda'''|$ are
at most $Cm$ on the disc. For real $\theta$, $\mathrm{Re}\,\Lambda(\theta)=\rho(\cos\theta)$, so
$\mathrm{Re}\,\Lambda'(\theta_0)=-\rho'(w_0)\sin\theta_0=0$ and
$\mathrm{Re}\,\Lambda''(\theta_0)=\rho''(w_0)\sin^2\theta_0\le-cm$. Put
\[
b=\mathrm{Im}\,\Lambda'(\theta_0),\qquad A=-\Lambda''(\theta_0),\qquad
T_{uv}=\frac{1}{\sqrt{2\pi A}}\exp\Bigl(\Lambda(\theta_0)-\frac{b^2}{2A}\Bigr),
\]
with the principal square root ($\mathrm{Re}\,A>0$).

\begin{lemma}\label{lem:laplace}
On $W_\lambda$: $|b|\le C\sqrt m$, $|T_{uv}|\ge c\,m^{-1/2}e^{\rho(w_0)}$, and
$t_{uv}=2\,\mathrm{Re}\,T_{uv}+O(|T_{uv}|\,m^{-1/5})$.
\end{lemma}

\begin{proof}
Write $a=u/m$, $a'=v/m$. Then $\Lambda'(\theta)/m$ differs by $O(1/m)$ from
$\ii z\bigl(\frac{a+a'}{z-1}+\frac{(1-a)P}{X}+\frac{(1-a')q}{Y}\bigr)-\ii$ at $z=e^{\ii\theta}$, which is a smooth
function of $(a,a',\theta)$ and vanishes at $a=a'=\frac14$, $\theta=\theta^*$: there $z^*=e^{\ii\theta^*}=\frac13+\frac{2\sqrt2}3\ii$, and
$\frac{1}{2(z^*-1)}+\frac{9}{16X(z^*)}+\frac{3}{16Y(z^*)}=\frac1{z^*}$, as one checks from $X(z^*)=\frac12+\frac{\sqrt2}2\ii$
and $Y(z^*)=\frac56+\frac{\sqrt2}6\ii$. Since $a-\frac14$, $a'-\frac14$ and $\theta_0-\theta^*$ are $O(m^{-1/2})$,
$|\Lambda'(\theta_0)|\le C\sqrt m$, which bounds $b$. Next, $\mathrm{Re}(b^2/A)=b^2\,\mathrm{Re}A/|A|^2\le C$, so
$|T_{uv}|\ge c\,|A|^{-1/2}|e^{\Lambda(\theta_0)}|\ge c\,m^{-1/2}e^{\rho(w_0)}$.

By Lemma~\ref{lem:integral}, $t_{uv}=\frac1{2\pi}\int_{-\pi}^\pi e^{\Lambda(\theta)}d\theta$, the integrand being
$Q_{uv}(e^{\ii\theta})e^{-\ii m\theta}$ on the whole circle. Split the interval into $J_+=\{|\theta-\theta_0|\le
m^{-2/5}\}$, $J_-=-J_+$ and the rest. On $J_+$, with $\theta=\theta_0+h$,
$\Lambda(\theta)=\Lambda(\theta_0)+\ii bh-\frac12Ah^2+R(h)$ with $|R(h)|\le Cm|h|^3\le Cm^{-1/5}$. Hence
\[
\frac1{2\pi}\int_{J_+}e^{\Lambda}=\frac{e^{\Lambda(\theta_0)}}{2\pi}\Bigl(\int_{\mathbb R}e^{\ii bh-Ah^2/2}dh
-\int_{|h|>m^{-2/5}}e^{\ii bh-Ah^2/2}dh\Bigr)+O\Bigl(e^{\rho(w_0)}m^{-1/5}\int_{\mathbb R}e^{-cmh^2}dh\Bigr).
\]
The first integral is $\sqrt{2\pi/A}\,e^{-b^2/(2A)}$, which makes the first term $T_{uv}$; the second is at most
$Ce^{-cm^{1/5}}$; the error term is $O(e^{\rho(w_0)}m^{-1/2-1/5})=O(|T_{uv}|m^{-1/5})$. The coefficients of $Q_{uv}$
are real, so the integral over $J_-$ is the complex conjugate. Off $J_+\cup J_-$, for $\theta\in[0,\pi]$,
$|\cos\theta-\cos\theta_0|=2\sin\frac{\theta+\theta_0}2\,\bigl|\sin\frac{\theta-\theta_0}2\bigr|\ge c|\theta-\theta_0|$
because $\theta_0$ stays away from $0$ and $\pi$, and likewise for $\theta\in[-\pi,0]$; since $\rho'(w_0)=0$ and
$\rho''\le-9n_2/128\le-cm$, the integrand there has modulus at most
$e^{\rho(w_0)-cm(\cos\theta-w_0)^2}\le e^{\rho(w_0)-cm^{1/5}}$, which is $O(|T_{uv}|m^{-1/5})$ after integration.
\end{proof}

\begin{lemma}\label{lem:phase}
Let $\alpha=\arccos\frac13$. If $(u,v)$ and $(u+1,v)$ lie in $W_\lambda$, then
$\arg T_{u+1,v}-\arg T_{uv}=\alpha+O(m^{-1/2})$ modulo $2\pi$.
\end{lemma}

\begin{proof}
Let $g(z)=(z-1)/X(z)$, so that $Q_{u+1,v}=gQ_{uv}$, and denote by $\Lambda^+$, $\theta_0^+$, $b^+$, $A^+$ the
quantities for $(u+1,v)$. Here $\Lambda^+=\Lambda+\ln g(e^{\ii\theta})$, and $w_0$ changes by $-4/(3m)$, so
$\theta_0^+-\theta_0=O(1/m)$. Hence
\[
\Lambda^+(\theta_0^+)-\Lambda(\theta_0)=\ln g(z_0)+\Lambda'(\theta_0)(\theta_0^+-\theta_0)+O(1/m)=\ln g(z_0)+O(m^{-1/2})
\]
by Lemma~\ref{lem:laplace}. Similarly $A^+-A=O(1)$ and $b^+-b=O(1)$, so $\arg A$ moves by $O(1/m)$ and $b^2/(2A)$ by
$O(|b|/m+b^2/m^2)=O(m^{-1/2})$. So $\arg T$ moves by $\mathrm{Im}\ln g(z_0)+O(m^{-1/2})$. Finally
$g(z^*)=\frac43z^*$, as $(z^*-1)\overline{X(z^*)}=\frac13+\frac{2\sqrt2}3\ii$ and $|X(z^*)|^2=\frac34$; thus
$\arg g(z^*)=\theta^*=\alpha$, and $z_0-z^*=O(m^{-1/2})$.
\end{proof}

The number $\alpha/\pi$ is irrational: by Niven's theorem~\cite{Niv}, if $\alpha/\pi$ is rational and $\cos\alpha$ is
rational then $\cos\alpha\in\{0,\pm\frac12,\pm1\}$.

\section{Proof of the theorem}\label{sec:proof}

The equality $\Inf_x(F_k)=(3/4)^{k-1}$ is Lemma~\ref{lem:sum}, and the upper bound is~\eqref{eq:upper}. For the lower
bound fix $\varepsilon\in(0,\frac14)$, $\lambda\ge1$ and an integer $L\ge1$. Since $\alpha/\pi$ is irrational, the
sequence $\ell\alpha/\pi$ is equidistributed modulo~$1$~\cite{Weyl16}, with discrepancy tending to $0$; the set of $\psi$
modulo $\pi$ with $|\cos\psi|<2\varepsilon$ is an interval of length $2\arcsin2\varepsilon\le2\pi\varepsilon$. So we may
take $L$ so large that for every real $\psi_0$ at most $3\varepsilon L$ of the numbers $\ell=0,\dots,L-1$ satisfy
$|\cos(\psi_0+\ell\alpha)|<2\varepsilon$.

Fix $v$ and cut the set of $u$ with $(u,v)\in W_\lambda$ into consecutive blocks of $L$ values, leaving at most $L-1$
values over. In a block starting at $u_0$, Lemma~\ref{lem:phase} gives
$\arg T_{u_0+\ell,v}=\arg T_{u_0,v}+\ell\alpha+O(Lm^{-1/2})$ for $\ell<L$, and for $m$ large the error is below
$\varepsilon$, so at most $3\varepsilon L$ values of $u$ in the block have $|\cos\arg T_{uv}|<\varepsilon$. On
$W_\lambda$, $\beta_{u+1}/\beta_u=(m-u)/(3(u+1))=1+O(m^{-1/2})$, so the weights $\beta_u$ in a block differ by a factor
at most $2$ for $m$ large, and the bad values carry at most $6\varepsilon$ of the block's weight. The left-over values
carry weight at most $(L-1)\max_u\beta_u=O(Lm^{-1/2})$ in each row. Summing over $v$, the pairs $(u,v)\in W_\lambda$ with
$|\cos\arg T_{uv}|<\varepsilon$, together with the left-over ones, have $\beta_u\beta_v$-weight at most
$6\varepsilon+o(1)$.

Let $(u,v)\in W_\lambda$ have $|\cos\arg T_{uv}|\ge\varepsilon$. By Lemma~\ref{lem:laplace},
$|t_{uv}|\ge2\varepsilon|T_{uv}|-C|T_{uv}|m^{-1/5}\ge\varepsilon|T_{uv}|\ge c\varepsilon m^{-1/2}e^{\rho(w_0)}$, and
$(u,v)\ne(0,0)$. With $2m-u-v=n_2-1$ and Lemma~\ref{lem:exponent},
\[
r_{uv}\ge\frac{2|t_{uv}|}{P^{n_2-1}}-\frac12\ge c\varepsilon m^{-1/2}\Bigl(\frac43\Bigr)^m-\frac12,
\]
which exceeds $1$ for $m$ large, and then $\Psi(r_{uv})\ge\log r_{uv}^2\ge2m\log\frac43-\log m-C_\varepsilon$. By
Chebyshev's inequality $\Pr[|u-m/4|>\lambda\sqrt m]\le\frac{3}{16\lambda^2}$, so $W_\lambda$ has weight at least
$(1-\frac3{16\lambda^2})^2$. Since $\Psi\ge0$, Lemma~\ref{lem:sum} gives
\[
\frac{E_x(F_k)}{\Inf_x(F_k)}\ge\Bigl(\bigl(1-\tfrac3{16\lambda^2}\bigr)^2-6\varepsilon-o(1)\Bigr)\Bigl(2m\log\frac43-\log
m-C_\varepsilon\Bigr).
\]
As $2m=k-1$, dividing by $(k-1)\log\frac43$ and letting $m\to\infty$, then $\lambda\to\infty$ and $\varepsilon\to0$,
gives $\liminf\ge1$. \qed

\begin{remark}
In $v$ the factor is $(z-1)/Y$, and $(z^*-1)/Y(z^*)=\frac43(-\frac13+\frac{2\sqrt2}3\ii)$, so the phase step is
$\pi-\alpha$ and the step from $(u,v)$ to $(u+1,v+1)$ is $\pi$: to first order the sign of $t_{uv}$ alternates along
diagonals.
\end{remark}

\section{Numerical checks}\label{sec:num}

Lemma~\ref{lem:sum} turns the ratio into a double sum of $(m+1)^2$ terms whose ingredients are integers over powers of
$4$, so the coefficients can be computed exactly in rational arithmetic; each term of $\Phi$ is then evaluated as
$a\log(1+b/a)$ with a $\log(1+\cdot)$ routine, since $a/b$ exceeds $10^{17}$ from $k=35$ on and is about $2\cdot10^{69}$ at $k=121$. The ratio $R_k=E_x(F_k)/\Inf_x(F_k)$ and the quotient
$R_k/((k-1)\log\frac43)$ are:
\begin{center}
\begin{tabular}{rrr@{\qquad}rrr}
$k$ & $R_k$ & quotient & $k$ & $R_k$ & quotient\\ \hline
$21$ & $6.631$ & $0.799$ & $101$ & $37.249$ & $0.898$\\
$41$ & $13.826$ & $0.833$ & $121$ & $45.276$ & $0.909$\\
$61$ & $21.448$ & $0.861$ & $133$ & $50.138$ & $0.915$\\
$81$ & $29.311$ & $0.883$ & $141$ & $53.319$ & $0.918$
\end{tabular}
\end{center}
For $k=3,5,7,9$ the values $1.0512$, $1.8769$, $2.5690$, $3.1179$ agree to $10^{-9}$ with the ratio computed from the
full Walsh--Hadamard transform of $F_k$ on $7$, $11$, $15$ and $19$ variables, which also confirms that $x$ is the
only variable in both a shortest minterm and a shortest maxterm. The Laplace approximation of
Lemma~\ref{lem:laplace} was compared with the exact $t_{uv}$ over all of $W_1$: the largest error relative to
$|T_{uv}|$ is $0.149$, $0.095$ and $0.071$ at $k=41$, $81$ and $141$, about $0.6\,m^{-1/2}$. Since $\rho$ is strictly concave and $\cos$ is decreasing on $(0,\pi)$,
$|Q_{uv}(e^{\ii\theta})|=e^{\rho(\cos\theta)}$ has no local maximum on $(0,\pi)$ other than $\theta_0$.

\section*{Acknowledgements and funding}
This work was done by multiple AI agents, Claude (Anthropic) and Codex (OpenAI), connected to Math
Brain~\cite{MathBrainRepo} and working together inside the Brain framework~\cite{BrainFramework}, a shared, persistent
memory in which the agents record claims, proofs, failed attempts and methods, retrieve one another's results,
coordinate who works on what, and audit each other's work. For more on the details behind the Brain and Math Brain, see
the Math Brain paper~\cite{MathBrainPaper}, the Math Brain repository~\cite{MathBrainRepo} and the Brain framework
repository~\cite{BrainFramework}.

The research presented in this paper was supported by the European Research Council (ERC) under the European Union's
Horizon Europe research and innovation programme (grant agreement No.~101041711), by the Simons Foundation as part of
the Collaboration on the Mathematical and Scientific Foundations of Deep Learning, by Heights Labs, by Convex Nexus
Capital, by the Israel Science Foundation (grant number 2258/19), and by the Israel Science Foundation (ISF Grant
4101/25). Computational resources were provided in part by Google Cloud through its education credits programme (credit
530550103).

\section*{Declaration of competing interest}
The author holds equity in Heights Labs and in Convex Nexus Capital and has a management role at Convex Nexus Capital.
These entities are included in the funding declaration above.

\section*{Declaration of generative AI and AI-assisted technologies}
The paper comes out of Math Brain, a project on the use of AI systems in research mathematics~\cite{BrainFramework,MathBrainRepo}. The family $F_{k,t}$
is from~\cite{Shm8}. Claude Opus 5.5 (Anthropic) found Theorem~\ref{thm:main} and its proof, wrote the numerical
checks and prepared the manuscript. The author directed the work and is responsible for the content of the paper.

\section*{Code and data availability}
Code and outputs: \url{https://github.com/yspennstate/fei-forced-coordinate}.

The paper is part of the Math Brain project, whose public repository is
\url{https://github.com/yspennstate/math-brain-public}~\cite{MathBrainRepo}.

\appendix
\section{Notation}\label{app:notation}

Symbols are listed roughly in the order in which they first appear, and the last column gives that place. The letters
$c$ and $g$ are reused, and each meaning has its own row.

\begin{longtable}{@{}>{\raggedright\arraybackslash}p{0.19\textwidth}>{\raggedright\arraybackslash}p{0.54\textwidth}>{\raggedright\arraybackslash}p{0.2\textwidth}@{}}
\toprule
Symbol & Meaning & First appears \\
\midrule
\endhead
$f$, $\hat f(S)$, $\chi_S$ & a Boolean function $\{0,1\}^n\to\{-1,1\}$, its Fourier coefficient at $S$, and the character at $S$ & Section~\ref{sec:intro} \\
$n$ & the number of variables of $f$ & Section~\ref{sec:intro} \\
$U$ & a random subset with $\Pr[U=S]=\hat f(S)^2$ & Section~\ref{sec:intro} \\
$C$ & an absolute constant, as in $H(U)\le C\,\E|U|$ & Section~\ref{sec:intro} \\
$x$ & a coordinate; for $F_k$, the coordinate lying in both a shortest minterm and a shortest maxterm & Section~\ref{sec:intro} \\
$E_x(f)$, $\Inf_x(f)$ & the erasure entropy $H(\mathbf 1[x\in U]\mid U\setminus\{x\})$, and the influence $\Pr[x\in U]$ & Section~\ref{sec:intro} \\
$h_2$ & the binary entropy & Section~\ref{sec:intro} \\
$F_{k,t}$ & the family of monotone functions of~\cite[Section~5]{Shm8} & Section~\ref{sec:intro} \\
$k$, $m$ & the odd number of inputs of the majority, and $m=(k-1)/2$ & Section~\ref{sec:intro} \\
$y^0$, $y^1$ & the point of $\{0,1\}^k$ with ones exactly in the first $m$ places, and $y^0+e_{m+1}$ & Section~\ref{sec:intro} \\
$\MAJ_k$, $M$ & the majority function on $k$ bits, and the same with $y^1$ removed from its true set & Section~\ref{sec:intro} \\
$\tilde f_k$, $F_k$ & $M(y)\vee(x\wedge[y\ge y^0])$, and the function on $2k+1$ variables obtained from it by replacing each $y_j$ with an OR or an AND of two fresh variables & Section~\ref{sec:intro} \\
$\Phi(a,b)$, $\Psi(r)$ & $a\log\frac{a+b}a+b\log\frac{a+b}b$, and $\Phi(r^2,1)$ & Section~\ref{sec:sum} \\
$P$, $q$ & $\frac34$ and $\frac14$, the probabilities that an OR and an AND of two uniform bits are true & Section~\ref{sec:sum} \\
$c(S)$ & the coefficients of $\tilde f_k$ in the biased basis & Section~\ref{sec:sum} \\
$\mu_j$, $\phi_j$, $\sigma$ & the bias of $y_j$, the biased basis function $(y-\mu_j)/\sigma$, and $\sigma=\sqrt{Pq}$ & Section~\ref{sec:sum} \\
$D$, $\bar A$, $N$ & the odd and even parts of $\tilde f_k$ in $x$, and the indicator $\mathbf 1[|y|\ge m+1]$ & Section~\ref{sec:sum} \\
$u$, $v$ & $|S\cap\{1,\dots,m\}|$ and $|S\cap\{m+2,\dots,k\}|$ & Section~\ref{sec:sum} \\
$\hat g(S)$ & the coefficient of a function $g$ at $\prod_{j\in S}\phi_j$ & Section~\ref{sec:sum} \\
$z$ & a complex variable & Section~\ref{sec:sum} \\
$X$, $Y$, $H_{uv}$ & $q+Pz$, $P+qz$, and $(z-1)^{u+v}X^{m-u}Y^{m+1-v}$ & Section~\ref{sec:sum} \\
$[z^c]$, $t_{uv}$ & the coefficient of $z^c$, and $\sum_{c\ge m+1}[z^c]H_{uv}$ & Lemma~\ref{lem:coef} \\
$\beta_u$, $r_{uv}$ & the binomial law $\binom mu q^uP^{m-u}$, and the quotient $|\hat{\bar A}(S)|/|\hat D(S)|$ for sets of type $(u,v)$ & Section~\ref{sec:sum} \\
$Q_{uv}$, $g_c$ & $H_{uv}/(z-1)$, and its coefficients, in the proof & Lemma~\ref{lem:integral} \\
$\theta$, $w$ & the angle with $z=e^{\ii\theta}$, and $w=\cos\theta$ & Lemma~\ref{lem:integral} \\
$n_1$, $n_2$ & $u+v-1$ and $2m+1-u-v$ & Section~\ref{sec:circle} \\
$\rho(w)$, $w_0$ & $\ln|Q_{uv}(e^{\ii\theta})|$ as a function of $w$, and its maximum point & Section~\ref{sec:circle} \\
$\lambda$, $W_\lambda$ & a width, and the pairs $(u,v)$ within $\lambda\sqrt m$ of $(m/4,m/4)$ in each coordinate & Section~\ref{sec:circle} \\
$c$, $C$ & positive constants that may depend on $\lambda$ but not on $m$ or $(u,v)$ & Section~\ref{sec:circle} \\
$\theta_0$, $\theta^*$, $z_0$ & $\arccos w_0$, $\arccos\frac13$, and $e^{\ii\theta_0}$ & Section~\ref{sec:circle} \\
$G(n_1)$ & the left side of the lemma as a function of $n_1$, in the proof & Lemma~\ref{lem:exponent} \\
$\Lambda(\theta)$, $\delta$ & $\ln Q_{uv}(e^{\ii\theta})-\ii m\theta$ on a small disc around $\theta^*$, and the radius of that disc & Section~\ref{sec:circle} \\
$b$, $A$, $T_{uv}$ & $\mathrm{Im}\,\Lambda'(\theta_0)$, $-\Lambda''(\theta_0)$, and the Laplace term $(2\pi A)^{-1/2}\exp(\Lambda(\theta_0)-b^2/(2A))$ & Section~\ref{sec:circle} \\
$a$, $a'$, $z^*$ & $u/m$, $v/m$, and $e^{\ii\theta^*}$, in the proof & Lemma~\ref{lem:laplace} \\
$J_\pm$, $h$, $R(h)$ & the arcs $|\theta\mp\theta_0|\le m^{-2/5}$, the offset $\theta-\theta_0$, and the cubic remainder, in the proof & Lemma~\ref{lem:laplace} \\
$\alpha$ & $\arccos\frac13$, an irrational multiple of $\pi$ & Lemma~\ref{lem:phase} \\
$g(z)$, $\Lambda^+$ & $(z-1)/X(z)$, and $\Lambda$ for the pair $(u+1,v)$, in the proof & Lemma~\ref{lem:phase} \\
$\varepsilon$, $L$ & a small number and a block length & Section~\ref{sec:proof} \\
$\psi_0$, $\ell$, $u_0$ & a phase, an index in a block, and the first $u$ of a block & Section~\ref{sec:proof} \\
$C_\varepsilon$ & a constant depending on $\varepsilon$ & Section~\ref{sec:proof} \\
$R_k$ & the ratio $E_x(F_k)/\Inf_x(F_k)$ & Section~\ref{sec:num} \\
\bottomrule
\end{longtable}

\begin{thebibliography}{99}
\bibitem{CKLS} S.~Chakraborty, R.~Kulkarni, S.~V.~Lokam, N.~Saurabh, Upper bounds on Fourier entropy, Theoret.
Comput. Sci. 654 (2016), 92--112.
\bibitem{FS09} P.~Flajolet, R.~Sedgewick, \emph{Analytic Combinatorics}, Cambridge University Press, 2009.
\bibitem{GMP} M.~J.~Gonz\'alez, P.~MacManus, M.~C.~Pereyra, Further evidence towards the Fourier entropy-influence
conjecture, arXiv:2606.00246, 2026.
\bibitem{FK96} E.~Friedgut, G.~Kalai, Every monotone graph property has a sharp threshold, Proc. Amer. Math. Soc.
124 (1996), 2993--3002.
\bibitem{Han} X.~Han, A new bound for the Fourier-entropy-influence conjecture, Combinatorica 45 (2025), no.~1,
Art.~4, doi:10.1007/s00493-024-00133-z; arXiv:2312.08271.
\bibitem{Hod} R.~Hod, Improved lower bounds for the Fourier entropy/influence conjecture via lexicographic functions,
arXiv:1711.00762, 2017.
\bibitem{KKLMS} E.~Kelman, G.~Kindler, N.~Lifshitz, D.~Minzer, M.~Safra, Towards a proof of the Fourier-entropy
conjecture?, Geom. Funct. Anal. 30 (2020), 1097--1138.
\bibitem{KMS12} N.~Keller, E.~Mossel, T.~Schlank, A note on the entropy/influence conjecture, Discrete Math. 312
(2012), 3364--3372, arXiv:1105.2651.
\bibitem{LH25} P.~Li, G.~Han, Strengthening Han's Fourier entropy-influence inequality via an
information-theoretic proof, arXiv:2512.03117, 2025.
\bibitem{MI26a} Mosaic Intelligence, A certified lower bound $C_{71}\ge 6.4901128435233943$ for the Fourier entropy-influence constant, Zenodo, 2026, doi:10.5281/zenodo.20794146.
\bibitem{MI26b} Mosaic Intelligence, A certified $n=17$ lower bound for the Fourier entropy-influence constant, Zenodo, 2026, doi:10.5281/zenodo.21113935.
\bibitem{Num26} Numaro, Improved lower bound $C_{71}>6.521845710923046575$ for the Fourier entropy-influence constant ($n=18$ explicit construction), certificate archive, Zenodo, 2026, doi:10.5281/zenodo.21497769.
\bibitem{OD14} R.~O'Donnell, \emph{Analysis of Boolean Functions}, Cambridge University Press, 2014.
\bibitem{OT13} R.~O'Donnell, L.-Y.~Tan, A composition theorem for the Fourier entropy-influence conjecture, ICALP 2013,
Lecture Notes in Comput. Sci. 7965, 780--791, doi:10.1007/978-3-642-39206-1\_66.
\bibitem{OWZ} R.~O'Donnell, J.~Wright, Y.~Zhou, The Fourier entropy-influence conjecture for certain classes of
Boolean functions, ICALP 2011, Lecture Notes in Comput. Sci. 6755, 330--341.
\bibitem{VW08} S.~Verd\'u, T.~Weissman, The information lost in erasures, IEEE Trans. Inform. Theory 54 (2008),
5030--5058.
\bibitem{WWW} A.~Wan, J.~Wright, C.~Wu, Decision trees, protocols, and the Fourier entropy-influence conjecture, ITCS 2014,
arXiv:1312.3003.
\bibitem{Weyl16} H.~Weyl, \"Uber die Gleichverteilung von Zahlen mod.~Eins, Math. Ann. 77 (1916), 313--352.
\bibitem{Niv} I.~Niven, \emph{Irrational Numbers}, Carus Mathematical Monographs 11, Mathematical Association of
America, 1956, Corollary~3.12.
\bibitem{Shm1} Y.~Shmalo, The one-coordinate step in the Fourier entropy-influence conjecture,
\url{https://github.com/yspennstate/fourier-entropy-influence-one-coordinate}, 2026.
\bibitem{Shm3} Y.~Shmalo, Good coordinates in the Fourier entropy-influence conjecture,
\url{https://github.com/yspennstate/fei-good-coordinates}, 2026.
\bibitem{Shm8} Y.~Shmalo, Hidden coordinates and selection rules in the Fourier entropy-influence conjecture,
\url{https://github.com/yspennstate/fei-hidden-coordinates}, 2026.
\bibitem{Shm10} Y.~Shmalo, The erasure entropy of a forced coordinate under every bias,
\url{https://github.com/yspennstate/fei-forced-coordinate-bias}, 2026.
\bibitem{ShmLTF} Y.~Shmalo, Fourier entropy of threshold functions and their compositions,
arXiv:2610.16069, 2026.
\bibitem{MathBrainRepo} Y.~Shmalo, \emph{Math Brain: public snapshot, code and tests}, 2026,
\url{https://github.com/yspennstate/math-brain-public}.
\bibitem{MathBrainPaper} Y.~Shmalo, \emph{Math Brain: shared, auditable memory for multi-agent mathematical research},
October 2026, \url{https://github.com/yspennstate/math-brain-public/tree/main/papers/math-brain-2026-10-24}.
\bibitem{BrainFramework} Y.~Shmalo, \emph{Brain Framework}, public software repository, 2026,
\url{https://github.com/yspennstate/brain-framework-public}.
\end{thebibliography}

\end{document}
